\section{Research priorities after the finite comparison}
\label{sec:outlook}
The strongest completed structural result is the complete-source coefficient
equation together with its commuting finite propagator. The forest polynomial
turns it into a quantitative scalar-filter interface. The corrected
Riemann--Siegel reference is the most concrete theta-specific input now
available on the microscopic clock. These results fit together, but their
composition still contains genuine acquisition requirements.

\begin{center}
\begin{tabular}{>{\raggedright\arraybackslash}p{0.28\linewidth}>{\raggedright\arraybackslash}p{0.61\linewidth}}
\toprule
Result & Assumption or remaining input \\
\midrule
Forest quotient & Polynomial algebra; a paid denominator margin for an enclosure.\\
Coefficient transport & Complete population and a common nonzero wall; integrated complete-source moments and errors.\\
Positive-time multiplicity transfer & Complete real adjacent-root motion; RH in the stated canonical global application.\\
Microscopic corrected reference & Imported real-time source theorem, fixed truncation chart; whole-wall margin still unacquired at growing height.\\
Finite theta certificate & Complete-source interval errors on the specified box; no extension to unbounded height.\\
\bottomrule
\end{tabular}
\end{center}

\subsection{A bounded next acquisition with a decisive question}
A useful next test would fix one genuinely high-height chart before examining
its zeros. For example, put \(\mathcal L=\log(x/(4\pi))\) and prescribe
\(t=\mathcal L^{-3}\), with spatial scales proportional to
\(\mathcal L^{-1}\). With natural unfolding
\(\nu\sim\mathcal L/(4\pi)\), this gives \(4t\nu^2\to0\).
The choice is a proposed experimental design, not a theorem about its
resulting roots. A fixed Riemann--Siegel cutoff, or explicitly paid cutoff
changes, is required.

The test should acquire the corrected normalized reference on every side of
the declared wall, its absolute source error, the complete population count,
coefficient errors, and the coupled forest denominator/residual margin.
It should report the resulting filter bound alongside an equally informed
direct fixed-time calculation. The decisive question is which term fails:
the wall margin, the source-to-coefficient conversion, degree growth, the
forest denominator, or the reference filter itself. One successful chart is
a validation of the acquisition method, not a limiting law. A broad sweep is
unjustified unless it tests a specified uniform estimate or cost model.

\subsection{The arithmetic estimate cannot be replaced by conditioning}
The long-term target is \eqref{eq:global-spectral-interface}: a vanishing
normalized filter and vanishing omitted mass on a family of complete
populations. Better source errors make a reference trustworthy; they do not
make its filter small. Even excellent whole-wall conditioning can enclose an
interior multiple zero. An independent arithmetic estimate must ultimately
control the relevant clustering of the complete reference, or the unchanged
field itself, on the prescribed joint height/time scale.

An averaged one-sided small-value estimate for the complete normalized field
and derivative is a possible direction, provided a proved consumer connects
it to the needed root counts and every interpolation, exceptional-set and
coverage loss is paid. This is a research suggestion, not an estimate supplied
here. Generic real-root heat dynamics, a mean zero count, or a statistical
model cannot stand in for the missing canonical input.

\subsection{Conditional roles for backward transport and curved paths}
A further backward calculation is justified when it can acquire the full
correlated residual in \eqref{eq:paid-duhamel}, including the explicit
reference heat defect, and improve the final certified error or cost relative
to direct acquisition from the same information. The finite propagator alone
does not establish that improvement. The scalar late-time implementation
with separately bounded algebraic errors has a different, unfavorable
amplification budget; its failure does not exclude a joint source estimate.

Complex time should be pursued only with a proposed complete-source error
bound that improves a specified path budget. Holomorphic continuation of the
reference is insufficient because the current imported field theorem is for
real time. A cap calculation must additionally retain every connector and
complement. The completed low-height range-quadrature pilot identifies a
spatial integration bottleneck, so further source precision or geometrical
rephrasing would not address its dominant error. Higher-order correlated
integration could be a bounded numerical study, with the same improvement
available to the direct baseline; it is a lower priority than the microscopic
arithmetic and conditioning questions.

The present stopping point is therefore a finite method with an explicit
source interface and an honest comparison, not an almost-complete proof of
simplicity. Density-one simplicity and every-zero simplicity remain distinct
targets; the latter cannot discard a sparse family of doubles. An RH-free
application must also establish endpoint reality independently.
